\begin{appendices}

\subsection{O*NET Descriptor Access Links}
\label{appendix:onet_urls}

Table~\ref{tab:onet_descriptor_urls} provides the access links and descriptor IDs for the six O*NET descriptors used in the three-dimensional feature mapping.

\begin{table}[H]
\centering
\caption{O*NET Descriptor Access Links and Descriptor IDs (for Reproducibility)}
\label{tab:onet_descriptor_urls}
\small
\begin{tabular}{@{}>{\raggedright\arraybackslash}p{2.2cm} >{\raggedright\arraybackslash}p{4.0cm} >{\raggedright\arraybackslash}p{2.3cm} >{\raggedright\arraybackslash}p{6.2cm}@{}}
\toprule
\textbf{Feature Axis} & \textbf{Descriptor Name} & \textbf{Descriptor ID} & \textbf{URL} \\
\midrule
$digit(x)$ & Working with Computers & 4.A.3.b.1 & \url{https://www.onetonline.org/find/descriptor/result/4.A.3.b.1} \\
 & Electronic Mail & 4.C.1.a.2.h & \url{https://www.onetonline.org/find/descriptor/result/4.C.1.a.2.h} \\
\addlinespace
$physical(x)$ & Using Hands to Handle Objects & 4.C.2.d.1.g & \url{https://www.onetonline.org/find/descriptor/result/4.C.2.d.1.g} \\
 & Kneeling/Crouching/Stooping & 4.C.2.d.1.e & \url{https://www.onetonline.org/find/descriptor/result/4.C.2.d.1.e} \\
\addlinespace
$iprisk(x)$ & Originality & 1.A.1.b.2 & \url{https://www.onetonline.org/find/descriptor/result/1.A.1.b.2} \\
 & Thinking Creatively & 4.A.2.b.2 & \url{https://www.onetonline.org/find/descriptor/result/4.A.2.b.2} \\
\bottomrule
\end{tabular}
\end{table}

\subsection{Two-Dimensional Projection Area Formula}
\label{appendix:area2d_formula}

To ensure the heterogeneity measure under two-dimensional projection is consistent with the three-dimensional definition, we define the two-dimensional triangle area for any projection $\pi$ (e.g., $(digit, physical)$, $(digit, iprisk)$, or $(physical, iprisk)$) as
\[
\mathrm{Area2D}_{\pi}(x_S,x_T,x_A)=\frac{1}{2}\left|\left(\pi(f(x_T))-\pi(f(x_S))\right)\times\left(\pi(f(x_A))-\pi(f(x_S))\right)\right|,
\]
where $\times$ denotes the cross product (scalar) operation of two-dimensional vectors, and $\pi(\cdot)$ is the projection mapping from three-dimensional features to two-dimensional coordinates.

\subsection{Supplementary Note: Task DNA Text Legends}
\label{app:task_dna_figures}

Due to page limitations, we did not list all task texts for the three occupations in the main text. When generating Figure~\ref{fig:task_dna_imfr_triptych}, we simultaneously exported weight-sorted task text legends for reproducibility and auditing purposes.

\subsection{Educational Program Recommendation Details}
\label{app:program_recommendations}

Table~\ref{tab:programs_genai_rationale} summarizes the recommended educational programs and selection rationale for the three representative occupations.

\begin{table}[H]
\centering
\renewcommand{\arraystretch}{1.3}
\small
\caption{Recommended Educational Programs and Selection Rationale}
\label{tab:programs_genai_rationale}
\begin{tabular}{@{}p{0.18\textwidth}p{0.26\textwidth}p{0.51\textwidth}@{}}
\toprule
\textbf{Occupation \& Institution Type} & \textbf{Recommended Institution \& Program} & \textbf{Selection Rationale} \\
\midrule

Robotics Engineers \newline(University) &
Carnegie Mellon University\newline
Robotics Institute: Master of Science in Robotic Systems Development (MRSD) &
\textbf{(1) Institution-level ``pathfinder'' attributes:} Since its establishment in 1979, the Robotics Institute has emphasized its leading position in robotics research and education~\cite{cmu_ri_2026}.\newline
\textbf{(2) Program objectives align with industry and systems engineering capabilities:} MRSD is positioned as an advanced industry-practice-oriented degree, emphasizing multidisciplinary technical and business capabilities, and includes internship pathways related to industry partners~\cite{cmu_mrsd_2026}.\newline
\textbf{(3) Gen-AI-compatible on-campus resources are explicitly available:} The university offers a ``Generative AI and Large Models'' graduate certificate program, which can serve as an institutionalized resource interface for introducing LLM/Gen-AI (e.g., natural language planning, human-robot collaboration, generative design) into robotics courses~\cite{cmu_gai_cert_2026}. \\
\addlinespace

First-Line Supervisors of Mechanics, Installers, and Repairers \newline(Vocational/Technical School) &
Universal Technical Institute (UTI)\newline
Automotive and Diesel Technology Programs &
\textbf{(1) Strong industry alignment and curriculum iteration capability:} UTI developed new BEV (Battery Electric Vehicle) courseware through collaboration with Bosch as part of its EV training strategy, demonstrating its ability to rapidly respond to industry technology transitions~\cite{uti_bosch_2022}.\newline
\textbf{(2) Realistic foundation for incorporating ``AI-enhanced diagnostics/damage estimation'' into skill contexts:} UTI's public content already directly discusses ``AI-supported diagnostics'' for identifying hidden damage and assisting in estimating repair parts and labor hours, indicating entry points for incorporating AI tool chains in educational narratives and training modules~\cite{uti_collision_ai_2026}. \\
\addlinespace

Court Reporters and Simultaneous Captioners \newline(Arts/Specialized School) &
Generations College (formerly MacCormac School)\newline
Court Reporting Associate Degree (A.A.S.) &
\textbf{(1) Strong program history and industry representation:} The school explicitly claims it established the ``first'' court reporting program of its kind in the United States in 1912, and emphasizes its industry reputation on enrollment pages~\cite{generations_court_2026}.\newline
\textbf{(2) Realtime training is explicitly required in program objectives:} The online program page explicitly states ``gain Realtime reporting experience in core courses''~\cite{generations_online_2026}.\newline
\textbf{(3) Connection points with industry quality standards exist:} The NCRA (National Court Reporters Association) published ``approved programs'' list includes Generations College (with approval validity period noted), serving as external endorsement evidence for program quality and realtime education commitment~\cite{ncra_approved_2026}. \\
\bottomrule
\end{tabular}
\end{table}

\end{appendices}
